% !TeX root = ../main.tex

\chapter{Conclusie}
\label{chap:conclusie}

Tijdens deze stage werd een werkende basisversie van het MSTC Race Engineer Dashboard
ontwikkeld. De applicatie leest live timingpagina's, normaliseert gegevens van verschillende
providers, bewaart een hervatbare racehistoriek en zet die data om in teamgerichte informatie. Het
dashboard toont onder meer recente rondes, sectoren, referentietijden, klassevergelijkingen,
pitstopinformatie, grafieken en automatische PDF-rapporten.

De belangrijkste technische bijdrage is de volledige keten van ruwe timingdata naar bruikbare
race-engineeringinformatie. Rondes en sectoren worden met context opgeslagen, niet-representatieve
rondes worden uitgesloten uit tempoanalyses en berekeningen zoals voorspellingen,
klassevergelijkingen en pitstopprojecties zitten in aparte modules. Daardoor bleef de software
testbaar en uitbreidbaar, ook toen de eisen tijdens de stage sterk evolueerden.

De test in Spa-Francorchamps bevestigde dat de dataketen met RIS Timing in een echte raceomgeving
bruikbaar was. Tegelijk kwamen daar precies de details naar boven die alleen live zichtbaar worden:
vlagtransities, pit-in- en outlaps, stabiele gaps en de selectie van geldige pace-ronden. De
opgeslagen Spa-data kon later gebruikt worden in regressietests, waardoor de applicatie ook na de race betrouwbaar bleef werken.

\pending{Na de \emph{24 Hours of Zolder}: voeg de definitieve eindconclusie toe. Vermeld of het
dashboard zijn einddoel behaalde, hoe het werkelijk werd gebruikt, welke beslissingen het
ondersteunde, welke beperkingen overbleven en welke feedback Jan en de piloten gaven.}

De stage vormde ook persoonlijk een mooie start in de motorsport. Volgend jaar werk ik aan mijn thesis bij \emph{Team WRT}, een van de grootste en meest succesvolle Belgische raceteams. Daar zal ik pitstops
analyseren met computer graphics en AI om de pitstops nog verder te optimaliseren. Deze
stage toonde al duidelijk waarom dat relevant is: ook in endurancewedstrijden, die de laatste jaren
steeds meer op lange sprintraces lijken, kan een trage pitstop een grote invloed hebben op de
uiteindelijke uitslag. Ik kijk er enorm naar uit om de kennis die ik in deze stage over motorsport heb opgedaan volgend jaar te gebruiken tijdens mijn thesis.
